\documentclass[11pt,a4paper]{article}
\usepackage[utf8]{inputenc}
\usepackage[T1]{fontenc}
\usepackage[ngerman]{babel}
\usepackage{helvet}
\renewcommand{\familydefault}{\sfdefault}
\usepackage{geometry}
\usepackage{hyperref}
\usepackage{booktabs}
\usepackage{longtable}
\usepackage{array}
\usepackage{enumitem}
\usepackage{xcolor}
\usepackage{fancyhdr}
\usepackage{titlesec}
\usepackage{tcolorbox}
\usepackage{tikz}
\usepackage{colortbl}
\usepackage{tabularx}
\usepackage{multirow}
\usepackage{amsmath}
\usepackage{amssymb}
\usepackage{siunitx}
\usepackage{pifont}

\usetikzlibrary{shapes.geometric, arrows.meta, positioning, calc}

\geometry{margin=1.8cm, top=2.2cm, bottom=1.8cm}
\hypersetup{colorlinks=true,linkcolor=physikblue,urlcolor=physikblue}

% Farben (Physik = Blau)
\definecolor{physikblue}{RGB}{44,90,160}
\definecolor{lightblue}{RGB}{220,235,250}
\definecolor{accentgreen}{RGB}{34,139,94}
\definecolor{accentorange}{RGB}{230,126,34}
\definecolor{lightgray}{RGB}{245,245,245}
\definecolor{positive}{RGB}{198,40,40}
\definecolor{negative}{RGB}{21,101,192}
\definecolor{examcolor}{RGB}{128,0,128}

% Symbole
\newcommand{\cmark}{\textcolor{accentgreen}{\ding{51}}}
\newcommand{\xmark}{\textcolor{red}{\ding{55}}}
\newcommand{\lmark}{\textcolor{accentorange}{\textbf{[LK]}}}
\newcommand{\checkbox}{$\square$}
\newcommand{\SE}{\textcolor{accentgreen}{\textbf{SE}}}
\newcommand{\DE}{\textcolor{physikblue}{\textbf{DE}}}

% Section Styling
\titleformat{\section}{\Large\bfseries\color{physikblue}}{\thesection.}{0.5em}{}[\vspace{-0.5em}\textcolor{physikblue}{\titlerule[1pt]}]
\titleformat{\subsection}{\large\bfseries\color{physikblue}}{\thesubsection}{0.5em}{}
\titleformat{\subsubsection}{\normalsize\bfseries\color{physikblue}}{\thesubsubsection}{0.5em}{}

% Header/Footer
\pagestyle{fancy}
\fancyhf{}
\fancyhead[L]{\small\color{gray}Leitfaden Physik MV}
\fancyhead[R]{\small\color{gray}\thepage}
\renewcommand{\headrulewidth}{0.4pt}

% Boxen
\newtcolorbox{infobox}[1][]{
  colback=lightblue, colframe=physikblue, fonttitle=\bfseries, title=#1,
  boxrule=0.5pt, arc=2pt, left=5pt, right=5pt, top=5pt, bottom=5pt
}
\newtcolorbox{wichtigbox}{
  colback=lightblue, colframe=physikblue, boxrule=0pt, leftrule=3pt, arc=0pt,
  left=8pt, right=5pt, top=5pt, bottom=5pt
}
\newtcolorbox{formelbox}[1][]{
  colback=lightgray, colframe=physikblue, fonttitle=\bfseries, title=#1,
  boxrule=1pt, arc=2pt, left=5pt, right=5pt, top=5pt, bottom=5pt
}
\newtcolorbox{experimentbox}{
  colback=accentgreen!10, colframe=accentgreen, boxrule=0.5pt, arc=2pt,
  left=5pt, right=5pt, top=3pt, bottom=3pt
}
\newtcolorbox{checklistbox}[1][]{
  colback=accentgreen!10, colframe=accentgreen, fonttitle=\bfseries, title=#1,
  boxrule=1pt, arc=2pt, left=5pt, right=5pt, top=5pt, bottom=5pt
}
\newtcolorbox{exambox}[1][]{
  colback=examcolor!10, colframe=examcolor, fonttitle=\bfseries, title=#1,
  boxrule=1pt, arc=2pt, left=5pt, right=5pt, top=5pt, bottom=5pt
}

\newcolumntype{H}{>{\columncolor{lightblue}\bfseries}l}
\newcolumntype{C}{>{\columncolor{lightblue}\bfseries}c}

% SI-Units Setup
\sisetup{locale=DE, per-mode=symbol}

\begin{document}

% === TITELSEITE ===
\begin{titlepage}
\centering
\vspace*{2.5cm}
{\Huge\bfseries\color{physikblue} Physik in\\[0.3em]Mecklenburg-Vorpommern\par}
\vspace{0.6cm}
{\LARGE\bfseries Umfassender Leitfaden\par}
\vspace{0.4cm}
{\large\textit{Sekundarstufe I und II mit Formeln, Experimenten und Prüfungsvorbereitung}\par}
\vspace{0.5cm}
{\large Rahmenpläne \textbullet{} Kompetenzen \textbullet{} Experimente \textbullet{} Abitur\par}
\vspace{1.5cm}

\begin{tcolorbox}[colback=lightgray, colframe=lightgray, width=13cm, boxrule=0pt]
\textbf{Stand:} Dezember 2025\\[0.2em]
\textbf{Grundlage:} Rahmenpläne Physik MV (Erprobungsfassung 2022/2023)\\[0.2em]
\textbf{Geltungsbereich:} Gymnasium, Gesamtschule, Regionale Schule\\[0.2em]
\textbf{Enthält:} Alle Formeln, Experimente, Operatoren, Prüfungshinweise
\end{tcolorbox}

\vspace{1.5cm}
\begin{tikzpicture}
\node[rectangle, rounded corners=5pt, fill=accentgreen!20, draw=accentgreen, line width=1pt, minimum width=10cm, minimum height=1.2cm] {
  \textbf{Merke:} $g \approx \SI{10}{\metre\per\second\squared}$ (kopfrechenbar!)
};
\end{tikzpicture}

\vfill
{\small Erstellt auf Grundlage der offiziellen Rahmenpläne des IQ M-V\par}
\end{titlepage}

% === INHALTSVERZEICHNIS ===
\tableofcontents
\newpage

% ============================================================================
% TEIL I: GRUNDLAGEN
% ============================================================================
\part*{Teil I: Übersicht und Struktur}
\addcontentsline{toc}{part}{Teil I: Übersicht und Struktur}

\section{Stundenverteilung}

\subsection{Sekundarstufe I (Klassen 7--10)}

\begin{center}
\begin{tabular}{c|c|l}
\rowcolor{lightblue}
\textbf{Klasse} & \textbf{USt} & \textbf{Themenfelder} \\
\hline
7 & $\sim$60 & Dichte, Kräfte, Energie, Kraftumformung \\
\rowcolor{lightgray}
8 & $\sim$60 & Licht, Elektrische Ladung, Stromkreise, Wärme \\
9 & $\sim$30 & Magnetismus, Induktion, Gleichförmige Bewegung \\
\rowcolor{lightgray}
10 & $\sim$60 & Beschleunigung, Impuls, Schwingungen, Wellen, Atomphysik \\
\hline
\multicolumn{3}{r}{\textbf{Gesamt: ca. 210 USt}} \\
\end{tabular}
\end{center}

\subsection{Sekundarstufe II (Qualifikationsphase)}

\begin{center}
\begin{tabular}{c|c|c|l}
\rowcolor{lightblue}
\textbf{Halbjahr} & \textbf{GK} & \textbf{LK} & \textbf{Themenfeld} \\
\hline
Q1 & 3 WSt & 5 WSt & Elektrisches und magnetisches Feld \\
\rowcolor{lightgray}
Q2 & 3 WSt & 5 WSt & Schwingungen und Wellen \\
Q3 & 3 WSt & 5 WSt & Quantenphysik \\
\rowcolor{lightgray}
Q4 & 3 WSt & 5 WSt & \lmark{} Spezielle Relativitätstheorie \\
\hline
\end{tabular}
\end{center}

\begin{wichtigbox}
\textbf{Wichtig:} Die Spezielle Relativitätstheorie (SRT) ist nur im \textbf{Leistungskurs} verbindlich!
\end{wichtigbox}

\section{Kompetenzbereiche}

\begin{center}
\begin{tikzpicture}[
  komp/.style={rectangle, rounded corners=3pt, minimum width=5cm, minimum height=1.2cm, text centered, font=\small},
]
\node[komp, fill=physikblue!25, draw=physikblue, line width=1pt] (s) at (0,0) {\textbf{[S] Sachkompetenz}\\Phänomene mit Modellen erklären};
\node[komp, fill=accentgreen!25, draw=accentgreen, line width=1pt, right=0.3cm of s] (e) {\textbf{[E] Erkenntnisgewinnung}\\Hypothesen, Experimente, Auswertung};
\node[komp, fill=accentorange!25, draw=accentorange, line width=1pt, below=0.4cm of s] (k) {\textbf{[K] Kommunikation}\\Fachsprache, Präsentation};
\node[komp, fill=examcolor!25, draw=examcolor, line width=1pt, right=0.3cm of k] (b) {\textbf{[B] Bewertung}\\Urteilen, Entscheiden, Reflektieren};
\end{tikzpicture}
\end{center}

% ============================================================================
% TEIL II: SEKUNDARSTUFE I
% ============================================================================
\newpage
\part*{Teil II: Sekundarstufe I -- Inhalte nach Klassen}
\addcontentsline{toc}{part}{Teil II: Sekundarstufe I}

\section{Klasse 7 ($\sim$60 USt)}

\subsection{Materie -- Dichte ($\sim$8 USt)}

\begin{formelbox}[Zentrale Formel]
\[
\rho = \frac{m}{V} \qquad \text{Dichte} = \frac{\text{Masse}}{\text{Volumen}}
\]
Einheit: $\si{\kilogram\per\metre\cubed}$ oder $\si{\gram\per\centi\metre\cubed}$
\end{formelbox}

\textbf{Verbindliche Inhalte:}
\begin{itemize}[leftmargin=0.5cm, itemsep=1pt]
\item Körper und Stoffe unterscheiden
\item Dichte als Proportionalitätsfaktor von Masse und Volumen
\item Sink-/Schweb-/Schwimmverhalten erklären
\end{itemize}

\begin{experimentbox}
\SE{} Masse und Volumen verschiedener Körper messen\\
\SE{} Materialuntersuchung durch Dichtebestimmung
\end{experimentbox}

\subsection{Wechselwirkung -- Kräfte ($\sim$18 USt)}

\begin{formelbox}[Zentrale Formeln]
\begin{align*}
F_G &= m \cdot g & \text{Gewichtskraft} \quad (g \approx \SI{10}{\newton\per\kilogram}) \\
F_S &= D \cdot s & \text{Federspannkraft (Hooke)} \\
F_R &= \mu \cdot F_N & \text{Reibungskraft} \\
p &= \frac{F}{A} & \text{Druck}
\end{align*}
\end{formelbox}

\textbf{Verbindliche Inhalte:}
\begin{itemize}[leftmargin=0.5cm, itemsep=1pt]
\item Kraft als Wechselwirkungsgröße (Betrag, Richtung, Angriffspunkt)
\item Gewichtskraft: $\SI{100}{\gram} \hat{=} \SI{1}{\newton}$
\item Haft-, Gleit-, Rollreibung
\item Auflage- und Tiefendruck, Luftdruck
\item Magnetische und elektrische Kräfte (qualitativ)
\end{itemize}

\begin{experimentbox}
\SE{} Wechselwirkungsprinzip nachweisen\\
\SE{} Federverlängerung messen (auch digital)\\
\SE{} Reibungskräfte vergleichen\\
\DE{} Druckabhängigkeit demonstrieren
\end{experimentbox}

\subsection{Energie -- Arbeit, Leistung, Wirkungsgrad ($\sim$20 USt)}

\begin{formelbox}[Zentrale Formeln]
\begin{align*}
E_{pot} &= m \cdot g \cdot h & \text{Lageenergie} \\
W &= F \cdot s & \text{Arbeit (bei $F \parallel s$)} \\
P &= \frac{W}{t} & \text{Leistung} \\
\eta &= \frac{E_{nutz}}{E_{zu}} & \text{Wirkungsgrad}
\end{align*}
\end{formelbox}

\textbf{Verbindliche Inhalte:}
\begin{itemize}[leftmargin=0.5cm, itemsep=1pt]
\item Energieformen: mechanisch, elektrisch, thermisch, chemisch, Kernenergie
\item Energieerhaltungssatz
\item Energieflussdiagramme
\item Goldene Regel der Mechanik
\end{itemize}

\subsection{Kraftumformende Einrichtungen ($\sim$8 USt)}

\textbf{Verbindliche Inhalte:}
\begin{itemize}[leftmargin=0.5cm, itemsep=1pt]
\item Geneigte Ebene
\item Ein- und zweiseitiger Hebel
\item Feste und lose Rolle, Flaschenzug
\end{itemize}

% === KLASSE 8 ===
\newpage
\section{Klasse 8 ($\sim$60 USt)}

\subsection{Wechselwirkung -- Licht ($\sim$10 USt)}

\textbf{Verbindliche Inhalte:}
\begin{itemize}[leftmargin=0.5cm, itemsep=1pt]
\item Spektrum und Farbwahrnehmung
\item Strahlenmodell des Lichts
\item Brechung und Dispersion
\item Bildentstehung am Auge
\item Lichterzeugung (Glühlampe, LED, Laser)
\end{itemize}

\begin{experimentbox}
\DE{} Kontinuierliches Spektrum erzeugen\\
\SE{} Strahlengang an Sammel-/Zerstreuungslinse\\
\SE{} Optisches Gerät bauen (Lupe, Fernrohr)
\end{experimentbox}

\subsection{Materie -- Elektrische Ladung ($\sim$10 USt)}

\textbf{Verbindliche Inhalte:}
\begin{itemize}[leftmargin=0.5cm, itemsep=1pt]
\item Elektrische Ladung als physikalische Größe
\item Coulombkraft (halbquantitativ)
\item Atomaufbau: Kern (Protonen, Neutronen), Hülle (Elektronen)
\item Ionen
\item Wirkungen und Gefahren des elektrischen Stroms
\end{itemize}

\subsection{System -- Stromkreise ($\sim$12 USt)}

\begin{formelbox}[Zentrale Formeln]
\begin{align*}
R &= \frac{U}{I} & \text{Ohm'sches Gesetz} \\
P_{el} &= U \cdot I & \text{Elektrische Leistung} \\
E_{el} &= P \cdot t & \text{Elektrische Energie}
\end{align*}
\textbf{Reihenschaltung:} $I = I_1 = I_2$, \quad $U_{ges} = U_1 + U_2$\\
\textbf{Parallelschaltung:} $U = U_1 = U_2$, \quad $I_{ges} = I_1 + I_2$
\end{formelbox}

\begin{experimentbox}
\SE{} Strom- und Spannungsmessung\\
\SE{} $I(U)$-Kennlinie aufnehmen
\end{experimentbox}

\subsection{Energie -- Temperatur und Wärme ($\sim$24 USt)}

\begin{formelbox}[Zentrale Formeln]
\begin{align*}
\Delta l &= \alpha \cdot l_0 \cdot \Delta T & \text{Längenänderung} \\
Q &= m \cdot c \cdot \Delta T & \text{Wärmemenge}
\end{align*}
\end{formelbox}

\textbf{Verbindliche Inhalte:}
\begin{itemize}[leftmargin=0.5cm, itemsep=1pt]
\item Temperatur und Kelvinskala
\item Volumenänderung bei Flüssigkeiten und Gasen
\item Anomalie des Wassers
\item Wärmeübertragung: Leitung, Strömung, Strahlung
\item Aggregatzustandsänderungen
\item Wärmepumpe, Verbrennungsmotor
\end{itemize}

% === KLASSE 9 ===
\newpage
\section{Klasse 9 ($\sim$30 USt)}

\subsection{Wechselwirkung -- Magnetismus ($\sim$8 USt)}

\textbf{Verbindliche Inhalte:}
\begin{itemize}[leftmargin=0.5cm, itemsep=1pt]
\item Elementarmagnete, Ferromagnetismus
\item Magnetisches Feld, Feldlinienmodell
\item Elektromagnetismus (Ørsted-Versuch)
\item Spule und gerader Leiter
\item Lorentzkraft (qualitativ)
\end{itemize}

\begin{experimentbox}
\SE{} Dipolteilung demonstrieren\\
\DE{} Ørsted-Versuch\\
\SE{} Elektromagnet: Abhängigkeiten untersuchen
\end{experimentbox}

\subsection{Energie -- Motor und Induktion ($\sim$8 USt)}

\textbf{Verbindliche Inhalte:}
\begin{itemize}[leftmargin=0.5cm, itemsep=1pt]
\item Gleichstrommotor (Aufbau und Funktion)
\item Elektromagnetische Induktion
\item Induktionsgesetz (qualitativ)
\item Generator und Transformator (Prinzip)
\end{itemize}

\subsection{Bewegungen -- Gleichförmig ($\sim$6 USt)}

\begin{formelbox}[Zentrale Formeln]
\begin{align*}
v &= \frac{\Delta s}{\Delta t} & \text{Geschwindigkeit} \\
s(t) &= v \cdot t + s_0 & \text{Weg-Zeit-Gesetz}
\end{align*}
\end{formelbox}

\textbf{Verbindliche Inhalte:}
\begin{itemize}[leftmargin=0.5cm, itemsep=1pt]
\item Bezugssysteme, Relativität der Bewegung
\item $s(t)$- und $v(t)$-Diagramme
\item Momentan- und Durchschnittsgeschwindigkeit
\end{itemize}

\subsection{Kontext -- Mit dem E-Bike unterwegs ($\sim$8 USt)}

\begin{formelbox}[Zusätzliche Formeln]
\begin{align*}
v &= \frac{2\pi r}{T} & \text{Bahngeschwindigkeit (Kreisbewegung)} \\
F_R &= \frac{m \cdot v^2}{r} & \text{Radialkraft}
\end{align*}
\end{formelbox}

% === KLASSE 10 ===
\newpage
\section{Klasse 10 ($\sim$60 USt)}

\subsection{Gleichmäßig beschleunigte Bewegung ($\sim$18 USt)}

\begin{formelbox}[Zentrale Formeln]
\begin{align*}
a &= \frac{\Delta v}{\Delta t} & \text{Beschleunigung} \\
v(t) &= a \cdot t + v_0 & \text{Geschwindigkeit-Zeit} \\
s(t) &= \frac{1}{2} a t^2 + v_0 t + s_0 & \text{Weg-Zeit} \\
v^2 &= 2 a s & \text{(ohne Zeit)}
\end{align*}
\textbf{Freier Fall:} $g = \SI{10}{\metre\per\second\squared}$
\end{formelbox}

\textbf{Verbindliche Inhalte:}
\begin{itemize}[leftmargin=0.5cm, itemsep=1pt]
\item Freier Fall
\item Waagerechter und schräger Wurf
\end{itemize}

\subsection{Impuls und Kraft ($\sim$12 USt)}

\begin{formelbox}[Zentrale Formeln]
\begin{align*}
p &= m \cdot v & \text{Impuls} \\
F &= \frac{\Delta p}{\Delta t} & \text{Kraft als Impulsänderung} \\
F &= m \cdot a & \text{Newton II}
\end{align*}
\textbf{Impulserhaltung:} $p_{vorher} = p_{nachher}$
\end{formelbox}

\textbf{Verbindliche Inhalte:}
\begin{itemize}[leftmargin=0.5cm, itemsep=1pt]
\item Trägheitssatz (Newton I)
\item Newton'sche Axiome
\item Impulserhaltung bei Stößen
\end{itemize}

\subsection{Mechanische Schwingungen ($\sim$12 USt)}

\begin{formelbox}[Zentrale Formeln]
\begin{align*}
T &= \frac{1}{f} & \text{Periode und Frequenz} \\
T &= 2\pi \sqrt{\frac{l}{g}} & \text{Fadenpendel} \\
T &= 2\pi \sqrt{\frac{m}{D}} & \text{Federpendel}
\end{align*}
\end{formelbox}

\textbf{Verbindliche Inhalte:} Amplitude, Periode, Frequenz, Eigenfrequenz, Resonanz

\subsection{Mechanische Wellen ($\sim$10 USt)}

\begin{formelbox}[Zentrale Formel]
\[
c = \lambda \cdot f \qquad \text{Ausbreitungsgeschwindigkeit}
\]
\end{formelbox}

\textbf{Verbindliche Inhalte:}
\begin{itemize}[leftmargin=0.5cm, itemsep=1pt]
\item Transversal- und Longitudinalwellen
\item Reflexion, Beugung, Interferenz
\end{itemize}

\subsection{Atom- und Kernphysik ($\sim$8 USt)}

\textbf{Verbindliche Inhalte:}
\begin{itemize}[leftmargin=0.5cm, itemsep=1pt]
\item Entwicklung der Atommodelle
\item Radioaktivität: $\alpha$-, $\beta$-, $\gamma$-Strahlung
\item Halbwertszeit
\item Kernspaltung und Kernfusion
\item Strahlenschutz
\end{itemize}

% ============================================================================
% TEIL III: SEKUNDARSTUFE II
% ============================================================================
\newpage
\part*{Teil III: Sekundarstufe II -- Qualifikationsphase}
\addcontentsline{toc}{part}{Teil III: Sekundarstufe II}

\section{Q1: Elektrisches und magnetisches Feld}

\begin{formelbox}[Zentrale Formeln Q1]
\begin{align*}
E &= \frac{F}{q} & \text{Elektrische Feldstärke} \\
C &= \frac{Q}{U} & \text{Kapazität} \\
E_C &= \frac{1}{2} C U^2 & \text{Energie im Kondensator} \\
F_L &= q \cdot v \cdot B & \text{Lorentzkraft} \\
U_{ind} &= -n \cdot \frac{d\Phi}{dt} & \text{Induktionsgesetz}
\end{align*}
\end{formelbox}

\textbf{Inhalte:}
\begin{itemize}[leftmargin=0.5cm, itemsep=1pt]
\item Homogenes elektrisches Feld (Plattenkondensator)
\item Potenzial und Spannung
\item Magnetische Feldstärke $B$
\item Bewegte Ladungen im Magnetfeld
\item Hall-Effekt
\item Selbstinduktion, Induktivität
\item Generator und Transformator (quantitativ)
\end{itemize}

\begin{experimentbox}
\textbf{Pflichtexperimente:}\\
Kondensatorauf-/-entladung \textbullet{} Elektronenstrahlablenkung \textbullet{} Fadenstrahlrohr
\end{experimentbox}

\section{Q2: Schwingungen und Wellen}

\begin{formelbox}[Zentrale Formeln Q2]
\begin{align*}
s(t) &= s_0 \cdot \sin(\omega t) & \text{Harmonische Schwingung} \\
\omega &= 2\pi f = \frac{2\pi}{T} & \text{Kreisfrequenz} \\
\lambda &= c \cdot T = \frac{c}{f} & \text{Wellenlänge} \\
d \cdot \sin\alpha &= n \cdot \lambda & \text{Beugung am Gitter}
\end{align*}
\end{formelbox}

\textbf{Inhalte:}
\begin{itemize}[leftmargin=0.5cm, itemsep=1pt]
\item Schwingungsgleichung, Energie bei Schwingungen
\item Gedämpfte und erzwungene Schwingung, Resonanz
\item Wellengleichung, Interferenz, stehende Wellen
\item Huygens'sches Prinzip, Beugung am Spalt/Gitter
\item Elektromagnetische Wellen: Hertz'scher Dipol, Spektrum, Polarisation
\end{itemize}

\section{Q3: Quantenphysik}

\begin{formelbox}[Zentrale Formeln Q3]
\begin{align*}
E &= h \cdot f & \text{Photonenenergie} \\
E_{kin} &= h \cdot f - W_A & \text{Photoeffekt} \\
\lambda &= \frac{h}{p} = \frac{h}{m \cdot v} & \text{De-Broglie-Wellenlänge} \\
N(t) &= N_0 \cdot e^{-\lambda t} & \text{Zerfallsgesetz} \\
E &= \Delta m \cdot c^2 & \text{Masse-Energie-Äquivalenz}
\end{align*}
\end{formelbox}

\textbf{Inhalte:}
\begin{itemize}[leftmargin=0.5cm, itemsep=1pt]
\item Welle-Teilchen-Dualismus
\item Photoeffekt, Photonen
\item Doppelspalt mit Elektronen
\item Bohr'sches Atommodell, Spektrallinien
\item Laser
\item Kernaufbau, Nuklidkarte, Bindungsenergie
\item Kernspaltung, Kernfusion, Reaktorsicherheit
\end{itemize}

\begin{experimentbox}
\textbf{Pflichtexperimente:}\\
Photoeffekt \textbullet{} Interferenz am Doppelspalt \textbullet{} Spektralanalyse
\end{experimentbox}

\section{Q4: Spezielle Relativitätstheorie \lmark{}}

\begin{formelbox}[Zentrale Formeln Q4 (nur LK)]
\begin{align*}
\gamma &= \frac{1}{\sqrt{1 - \frac{v^2}{c^2}}} & \text{Lorentzfaktor} \\
\Delta t' &= \gamma \cdot \Delta t & \text{Zeitdilatation} \\
l' &= \frac{l}{\gamma} & \text{Längenkontraktion} \\
E &= m \cdot c^2 & \text{Ruheenergie} \\
E_{ges} &= \gamma \cdot m \cdot c^2 & \text{Gesamtenergie}
\end{align*}
\end{formelbox}

\textbf{Inhalte:}
\begin{itemize}[leftmargin=0.5cm, itemsep=1pt]
\item Postulate der SRT
\item Zeitdilatation und Längenkontraktion
\item Relativistische Masse und Energie
\end{itemize}

% ============================================================================
% TEIL IV: OPERATOREN UND BEWERTUNG
% ============================================================================
\newpage
\part*{Teil IV: Operatoren und Bewertung}
\addcontentsline{toc}{part}{Teil IV: Operatoren und Bewertung}

\section{Anforderungsbereiche (AFB)}

\begin{center}
\begin{tikzpicture}[
  block/.style={rectangle, rounded corners=3pt, minimum width=3.5cm, minimum height=1.2cm, text centered},
]
\node[block, fill=accentgreen!25, draw=accentgreen, line width=1pt] (afb1) at (0,0) {\textbf{AFB I}\\Reproduktion\\$\sim$30\%};
\node[block, fill=accentorange!25, draw=accentorange, line width=1pt, right=0.3cm of afb1] (afb2) {\textbf{AFB II}\\Reorganisation\\$\sim$40\%};
\node[block, fill=physikblue!25, draw=physikblue, line width=1pt, right=0.3cm of afb2] (afb3) {\textbf{AFB III}\\Transfer\\$\sim$30\%};

\node[below=0.15cm of afb1, font=\footnotesize, text width=3.5cm, align=center] {nennen, beschreiben,\\angeben, berechnen};
\node[below=0.15cm of afb2, font=\footnotesize, text width=3.5cm, align=center] {erklären, ermitteln,\\vergleichen, auswerten};
\node[below=0.15cm of afb3, font=\footnotesize, text width=3.5cm, align=center] {analysieren, begründen,\\beurteilen, entwickeln};
\end{tikzpicture}
\end{center}

\section{Operatoren}

\begin{center}
\footnotesize
\begin{tabularx}{\textwidth}{c|l|X}
\rowcolor{lightblue}
\textbf{AFB} & \textbf{Operator} & \textbf{Definition} \\
\hline
I & nennen & Aufzählen ohne Erläuterung \\
\rowcolor{lightgray}
I & beschreiben & Sachverhalt in eigenen Worten wiedergeben \\
I & angeben & Ergebnis ohne Lösungsweg nennen \\
\rowcolor{lightgray}
I & berechnen & Ergebnis rechnerisch mit Lösungsweg ermitteln \\
\hline
II & erklären & Sachverhalt verständlich und nachvollziehbar machen \\
\rowcolor{lightgray}
II & ermitteln & Ergebnis aus Daten, Diagrammen oder Texten gewinnen \\
II & bestimmen & Lösung durch Anwenden bekannter Verfahren herleiten \\
\rowcolor{lightgray}
II & vergleichen & Gemeinsamkeiten und Unterschiede herausarbeiten \\
II & auswerten & Daten oder Messwerte interpretieren \\
\hline
\rowcolor{lightgray}
III & analysieren & Systematisch und gezielt untersuchen \\
III & begründen & Aussage durch physikalische Argumente stützen \\
\rowcolor{lightgray}
III & beurteilen & Sachverhalt nach Kriterien einschätzen und bewerten \\
III & entwickeln & Lösungsstrategie selbstständig erarbeiten \\
\rowcolor{lightgray}
III & diskutieren & Pro- und Kontra-Argumente gegenüberstellen und abwägen \\
\hline
\end{tabularx}
\end{center}

\section{Bewertung (15-NP-Skala Sek II)}

\begin{center}
\footnotesize
\begin{tabular}{c|c|c|c|c|c|c|c|c|c|c|c|c|c|c|c|c}
\rowcolor{lightblue}
\textbf{NP} & 15 & 14 & 13 & 12 & 11 & 10 & 9 & 8 & 7 & 6 & 5 & 4 & 3 & 2 & 1 & 0 \\
\hline
\textbf{\%} & $\geq$99 & $\geq$97 & $\geq$96 & $\geq$91 & $\geq$85 & $\geq$80 & $\geq$73 & $\geq$67 & $\geq$60 & $\geq$53 & $\geq$47 & $\geq$40 & $\geq$33 & $\geq$27 & $\geq$20 & <20 \\
\rowcolor{lightgray}
\textbf{Note} & 1+ & 1 & 1- & 2+ & 2 & 2- & 3+ & 3 & 3- & 4+ & 4 & 4- & 5+ & 5 & 5- & 6 \\
\hline
\end{tabular}
\end{center}

% ============================================================================
% TEIL V: KONSTANTEN UND FORMELN
% ============================================================================
\newpage
\part*{Teil V: Wichtige Konstanten}
\addcontentsline{toc}{part}{Teil V: Konstanten und Formeln}

\section{Physikalische Konstanten}

\begin{center}
\begin{tabular}{l|c|l}
\rowcolor{lightblue}
\textbf{Konstante} & \textbf{Symbol} & \textbf{Wert} \\
\hline
Fallbeschleunigung (Erde) & $g$ & $\approx \SI{10}{\metre\per\second\squared}$ (kopfrechenbar!) \\
\rowcolor{lightgray}
Lichtgeschwindigkeit & $c$ & $\SI{3,0e8}{\metre\per\second}$ \\
Elementarladung & $e$ & $\SI{1,6e-19}{\coulomb}$ \\
\rowcolor{lightgray}
Planck-Konstante & $h$ & $\SI{6,63e-34}{\joule\second}$ \\
Elektronenmasse & $m_e$ & $\SI{9,1e-31}{\kilogram}$ \\
\rowcolor{lightgray}
Protonenmasse & $m_p$ & $\SI{1,67e-27}{\kilogram}$ \\
Gravitationskonstante & $G$ & $\SI{6,67e-11}{\newton\metre\squared\per\kilogram\squared}$ \\
\rowcolor{lightgray}
Elektrische Feldkonstante & $\varepsilon_0$ & $\SI{8,85e-12}{\ampere\second\per\volt\per\metre}$ \\
Magnetische Feldkonstante & $\mu_0$ & $4\pi \times 10^{-7}\,\si{\volt\second\per\ampere\per\metre}$ \\
\hline
\end{tabular}
\end{center}

\section{Kopfrechenbare Werte}

\begin{wichtigbox}
\textbf{IMMER verwenden für Aufgaben:}
\begin{itemize}[leftmargin=0.5cm, itemsep=2pt]
\item $g = \SI{10}{\metre\per\second\squared}$ (NICHT 9,81!)
\item Geschwindigkeiten: 36, 72, 90, 108 km/h $\rightarrow$ 10, 20, 25, 30 m/s
\item Zeiten: 2, 3, 4, 5, 6, 8, 10 s
\item Beschleunigungen: 2, 4, 5, 10 m/s²
\item Wurzeln: $\sqrt{4}=2$, $\sqrt{9}=3$, $\sqrt{16}=4$, $\sqrt{25}=5$, $\sqrt{36}=6$, $\sqrt{100}=10$
\end{itemize}
\end{wichtigbox}

% ============================================================================
% TEIL VI: PRÜFUNGSHINWEISE
% ============================================================================
\newpage
\part*{Teil VI: Examenshinweise}
\addcontentsline{toc}{part}{Teil VI: Examenshinweise}

\begin{exambox}[Typische mündliche Prüfungsfragen (2. Staatsexamen)]
\begin{enumerate}[leftmargin=0.5cm, itemsep=3pt]
\item „Erläutern Sie den Aufbau des Rahmenplans Physik Sek I."
\item „Wie führen Sie das Konzept der Energie in Klasse 7 ein?"
\item „Beschreiben Sie ein Schülerexperiment zur Dichtebestimmung."
\item „Wie differenzieren Sie zwischen GK und LK in der Quantenphysik?"
\item „Nennen Sie drei AFB-III-Operatoren mit Beispielaufgaben."
\item „Wie setzen Sie digitale Messwerterfassung im Unterricht ein?"
\end{enumerate}
\end{exambox}

\section{Beispielantworten}

\textbf{Frage 2: Energiekonzept Klasse 7}

„Ich führe Energie als \textit{Erhaltungsgröße} ein, nicht als abstrakte Definition. Konkret:
\begin{itemize}[leftmargin=0.5cm, itemsep=1pt]
\item \textbf{Einstieg:} Alltagsbeispiele (Fahrrad bergab, Gummiband)
\item \textbf{Energieformen:} Erst qualitativ benennen (Bewegungs-, Lage-, Wärmeenergie)
\item \textbf{Umwandlung:} Energieflussdiagramme zeichnen
\item \textbf{Erhaltung:} An geschlossenem System demonstrieren
\item \textbf{Quantifizierung:} Erst danach $E_{pot} = mgh$ einführen
\end{itemize}
Die Formel $W = F \cdot s$ kommt erst nach dem Energiekonzept, da Arbeit als Energieübertragung verstanden werden soll."

\textbf{Frage 3: Schülerexperiment Dichte}

„Die SuS bestimmen die Dichte verschiedener Metallwürfel:
\begin{enumerate}[leftmargin=0.5cm, itemsep=1pt]
\item Masse mit Waage messen
\item Volumen über Kantenlänge oder Überlaufmethode bestimmen
\item $\rho = m/V$ berechnen
\item Mit Tabellenwerten vergleichen $\rightarrow$ Material identifizieren
\end{enumerate}
\textbf{Kompetenzen:} [E] Messen und Auswerten, [K] Protokollieren, [S] Dichte anwenden"

% ============================================================================
% TEIL VII: CHECKLISTEN
% ============================================================================
\newpage
\part*{Teil VII: Checklisten}
\addcontentsline{toc}{part}{Teil VII: Checklisten}

\begin{checklistbox}[Checkliste: Vor dem Unterricht (Sek I)]
\begin{itemize}[leftmargin=0.5cm, label=\checkbox]
\item Themenfeld und USt-Umfang im RP geprüft?
\item Verbindliche Inhalte identifiziert?
\item Pflichtexperimente (SE/DE) eingeplant?
\item Kopfrechenbare Zahlenwerte gewählt ($g = 10$)?
\item Kompetenzen [S], [E], [K], [B] berücksichtigt?
\item Querschnittsthemen integriert (BNE, PG, MD)?
\end{itemize}
\end{checklistbox}

\vspace{0.5cm}

\begin{checklistbox}[Checkliste: Vor dem Unterricht (Sek II)]
\begin{itemize}[leftmargin=0.5cm, label=\checkbox]
\item Halbjahr (Q1--Q4) und Themenfeld korrekt?
\item GK/LK-Unterschiede beachtet (SRT nur LK)?
\item Operatoren aus verbindlicher Liste verwendet?
\item Pflichtexperimente eingeplant?
\item Vorabhinweise für Abitur geprüft?
\item AFB-Verteilung 30/40/30 beachtet?
\end{itemize}
\end{checklistbox}

\vspace{0.5cm}

\begin{checklistbox}[Checkliste: Klausurerstellung]
\begin{itemize}[leftmargin=0.5cm, label=\checkbox]
\item AFB-Verteilung eingehalten ($\sim$30/40/30)?
\item Operatoren korrekt verwendet?
\item Kopfrechenbare Werte ($g = 10$, einfache Wurzeln)?
\item Einheiten bei allen Angaben?
\item Erwartungshorizont mit Punktverteilung?
\item Lösungsweg erforderlich bei „berechnen"?
\end{itemize}
\end{checklistbox}

% ============================================================================
% THEMEN-CHECKLISTE
% ============================================================================
\newpage
\section{Themen-Checkliste: Soll und Optional}

\subsection{Sekundarstufe I}

\begin{center}
\textbf{Klasse 7} (ca. 60 USt)
\end{center}

\begin{tabularx}{\textwidth}{c|l|c|X}
\rowcolor{lightblue}
\checkbox & \textbf{Thema} & \textbf{USt} & \textbf{Status} \\
\hline
\checkbox & Materie -- Dichte & $\sim$8 & \textbf{SOLL} \\
\rowcolor{lightgray}
\checkbox & Wechselwirkung -- Kräfte & $\sim$18 & \textbf{SOLL} \\
\checkbox & Energie -- Arbeit, Leistung, Wirkungsgrad & $\sim$20 & \textbf{SOLL} \\
\rowcolor{lightgray}
\checkbox & Kraftumformende Einrichtungen & $\sim$8 & \textbf{SOLL} \\
\checkbox & \textit{Vertiefung: Hydrostatik} & -- & \textit{optional} \\
\hline
\end{tabularx}

\vspace{0.4cm}

\begin{center}
\textbf{Klasse 8} (ca. 60 USt)
\end{center}

\begin{tabularx}{\textwidth}{c|l|c|X}
\rowcolor{lightblue}
\checkbox & \textbf{Thema} & \textbf{USt} & \textbf{Status} \\
\hline
\checkbox & Wechselwirkung -- Licht & $\sim$10 & \textbf{SOLL} \\
\rowcolor{lightgray}
\checkbox & Materie -- Elektrische Ladung & $\sim$10 & \textbf{SOLL} \\
\checkbox & System -- Stromkreise & $\sim$12 & \textbf{SOLL} \\
\rowcolor{lightgray}
\checkbox & Energie -- Temperatur und Wärme & $\sim$24 & \textbf{SOLL} \\
\checkbox & \textit{Vertiefung: Optische Geräte bauen} & -- & \textit{optional} \\
\rowcolor{lightgray}
\checkbox & \textit{Vertiefung: Halbleiter} & -- & \textit{optional} \\
\hline
\end{tabularx}

\vspace{0.4cm}

\begin{center}
\textbf{Klasse 9} (ca. 30 USt)
\end{center}

\begin{tabularx}{\textwidth}{c|l|c|X}
\rowcolor{lightblue}
\checkbox & \textbf{Thema} & \textbf{USt} & \textbf{Status} \\
\hline
\checkbox & Wechselwirkung -- Magnetismus & $\sim$8 & \textbf{SOLL} \\
\rowcolor{lightgray}
\checkbox & Energie -- Motor und Induktion & $\sim$8 & \textbf{SOLL} \\
\checkbox & Bewegungen -- Gleichförmig & $\sim$6 & \textbf{SOLL} \\
\rowcolor{lightgray}
\checkbox & Kontext: Mit dem E-Bike unterwegs & $\sim$8 & \textbf{SOLL} \\
\checkbox & \textit{Vertiefung: Wechselstrom} & -- & \textit{optional} \\
\hline
\end{tabularx}

\vspace{0.4cm}

\begin{center}
\textbf{Klasse 10} (ca. 60 USt)
\end{center}

\begin{tabularx}{\textwidth}{c|l|c|X}
\rowcolor{lightblue}
\checkbox & \textbf{Thema} & \textbf{USt} & \textbf{Status} \\
\hline
\checkbox & Gleichmäßig beschleunigte Bewegung & $\sim$18 & \textbf{SOLL} \\
\rowcolor{lightgray}
\checkbox & Impuls und Kraft (Newton) & $\sim$12 & \textbf{SOLL} \\
\checkbox & Mechanische Schwingungen & $\sim$12 & \textbf{SOLL} \\
\rowcolor{lightgray}
\checkbox & Mechanische Wellen & $\sim$10 & \textbf{SOLL} \\
\checkbox & Atom- und Kernphysik & $\sim$8 & \textbf{SOLL} \\
\rowcolor{lightgray}
\checkbox & \textit{Vertiefung: Schräger Wurf (quantitativ)} & -- & \textit{optional} \\
\checkbox & \textit{Vertiefung: Akustik} & -- & \textit{optional} \\
\hline
\end{tabularx}

\subsection{Sekundarstufe II (Qualifikationsphase)}

\begin{center}
\textbf{Q1: Elektrisches und magnetisches Feld}
\end{center}

\begin{tabularx}{\textwidth}{c|l|c|X}
\rowcolor{lightblue}
\checkbox & \textbf{Thema} & \textbf{Kurs} & \textbf{Status} \\
\hline
\checkbox & Homogenes elektrisches Feld & GK + LK & \textbf{SOLL} \\
\rowcolor{lightgray}
\checkbox & Potenzial und Spannung & GK + LK & \textbf{SOLL} \\
\checkbox & Kondensator (Auf-/Entladung) & GK + LK & \textbf{SOLL} \\
\rowcolor{lightgray}
\checkbox & Magnetische Feldstärke B & GK + LK & \textbf{SOLL} \\
\checkbox & Lorentzkraft, Elektronenstrahlablenkung & GK + LK & \textbf{SOLL} \\
\rowcolor{lightgray}
\checkbox & Hall-Effekt & GK + LK & \textbf{SOLL} \\
\checkbox & Fadenstrahlrohr (e/m-Bestimmung) & GK + LK & \textbf{SOLL} \\
\rowcolor{lightgray}
\checkbox & Induktion, Induktionsgesetz & GK + LK & \textbf{SOLL} \\
\checkbox & Generator, Transformator (quantitativ) & GK + LK & \textbf{SOLL} \\
\rowcolor{lightgray}
\checkbox & \textit{Wirbelströme (Vertiefung)} & LK & \textit{optional} \\
\checkbox & \textit{Massenspektrometer} & LK & \textit{optional} \\
\hline
\end{tabularx}

\vspace{0.4cm}

\begin{center}
\textbf{Q2: Schwingungen und Wellen}
\end{center}

\begin{tabularx}{\textwidth}{c|l|c|X}
\rowcolor{lightblue}
\checkbox & \textbf{Thema} & \textbf{Kurs} & \textbf{Status} \\
\hline
\checkbox & Harmonische Schwingung (s, v, a) & GK + LK & \textbf{SOLL} \\
\rowcolor{lightgray}
\checkbox & Schwingungsenergie & GK + LK & \textbf{SOLL} \\
\checkbox & Gedämpfte/erzwungene Schwingung, Resonanz & GK + LK & \textbf{SOLL} \\
\rowcolor{lightgray}
\checkbox & Wellengleichung, Interferenz & GK + LK & \textbf{SOLL} \\
\checkbox & Stehende Wellen & GK + LK & \textbf{SOLL} \\
\rowcolor{lightgray}
\checkbox & Huygens'sches Prinzip & GK + LK & \textbf{SOLL} \\
\checkbox & Beugung am Spalt und Gitter & GK + LK & \textbf{SOLL} \\
\rowcolor{lightgray}
\checkbox & EM-Wellen, Hertz'scher Dipol & GK + LK & \textbf{SOLL} \\
\checkbox & Polarisation & GK + LK & \textbf{SOLL} \\
\rowcolor{lightgray}
\checkbox & \textit{Gekoppelte Schwingungen} & LK & \textit{optional} \\
\checkbox & \textit{Dopplereffekt} & LK & \textit{optional} \\
\hline
\end{tabularx}

\vspace{0.4cm}

\begin{center}
\textbf{Q3: Quantenphysik}
\end{center}

\begin{tabularx}{\textwidth}{c|l|c|X}
\rowcolor{lightblue}
\checkbox & \textbf{Thema} & \textbf{Kurs} & \textbf{Status} \\
\hline
\checkbox & Photoeffekt, Photonen & GK + LK & \textbf{SOLL} \\
\rowcolor{lightgray}
\checkbox & h-Bestimmung & GK + LK & \textbf{SOLL} \\
\checkbox & Welle-Teilchen-Dualismus & GK + LK & \textbf{SOLL} \\
\rowcolor{lightgray}
\checkbox & De-Broglie-Wellenlänge & GK + LK & \textbf{SOLL} \\
\checkbox & Doppelspalt mit Elektronen & GK + LK & \textbf{SOLL} \\
\rowcolor{lightgray}
\checkbox & Bohr'sches Atommodell, Spektrallinien & GK + LK & \textbf{SOLL} \\
\checkbox & Laser (Prinzip) & GK + LK & \textbf{SOLL} \\
\rowcolor{lightgray}
\checkbox & Kernaufbau, Nuklidkarte & GK + LK & \textbf{SOLL} \\
\checkbox & Bindungsenergie, Massendefekt & GK + LK & \textbf{SOLL} \\
\rowcolor{lightgray}
\checkbox & Zerfallsgesetz, Halbwertszeit & GK + LK & \textbf{SOLL} \\
\checkbox & Kernspaltung, Kernfusion & GK + LK & \textbf{SOLL} \\
\rowcolor{lightgray}
\checkbox & \textit{Heisenberg'sche Unschärferelation} & LK & \textit{optional} \\
\checkbox & \textit{Franck-Hertz-Versuch} & LK & \textit{optional} \\
\hline
\end{tabularx}

\vspace{0.4cm}

\begin{center}
\textbf{Q4: Spezielle Relativitätstheorie} \lmark{}
\end{center}

\begin{tabularx}{\textwidth}{c|l|c|X}
\rowcolor{lightblue}
\checkbox & \textbf{Thema} & \textbf{Kurs} & \textbf{Status} \\
\hline
\checkbox & Postulate der SRT & \textbf{nur LK} & \textbf{SOLL (LK)} \\
\rowcolor{lightgray}
\checkbox & Zeitdilatation & \textbf{nur LK} & \textbf{SOLL (LK)} \\
\checkbox & Längenkontraktion & \textbf{nur LK} & \textbf{SOLL (LK)} \\
\rowcolor{lightgray}
\checkbox & Relativistische Masse & \textbf{nur LK} & \textbf{SOLL (LK)} \\
\checkbox & $E = mc^2$, Ruheenergie & \textbf{nur LK} & \textbf{SOLL (LK)} \\
\rowcolor{lightgray}
\checkbox & \textit{Myonenzerfall (Anwendung)} & \textbf{nur LK} & \textit{optional} \\
\checkbox & \textit{Lorentz-Transformation} & \textbf{nur LK} & \textit{optional} \\
\hline
\end{tabularx}

\vspace{0.5cm}

\begin{wichtigbox}
\textbf{Legende:}
\begin{itemize}[leftmargin=0.5cm, itemsep=1pt]
\item \textbf{SOLL} = Verbindlicher Inhalt laut Rahmenplan (muss unterrichtet werden)
\item \textit{optional} = Vertiefung/Erweiterung (kann unterrichtet werden, wenn Zeit vorhanden)
\item \lmark{} = Nur Leistungskurs (Q4 ist im GK nicht verbindlich)
\end{itemize}
\end{wichtigbox}

% ============================================================================
% QUICK-REFERENCE
% ============================================================================
\newpage
\section*{QUICK-REFERENCE: Physik MV auf einen Blick}
\addcontentsline{toc}{section}{Quick-Reference (1 Seite)}

\begin{center}
\fbox{\parbox{0.95\textwidth}{
\centering\Large\bfseries\color{physikblue} PHYSIK MECKLENBURG-VORPOMMERN -- QUICK-REFERENCE
}}
\end{center}

\vspace{0.3cm}

\begin{minipage}[t]{0.48\textwidth}
\textbf{\color{physikblue}SEKUNDARSTUFE I}
\begin{itemize}[leftmargin=0.3cm, itemsep=1pt, topsep=2pt]
\item \textbf{Kl. 7 (60 USt):} Dichte, Kräfte, Energie
\item \textbf{Kl. 8 (60 USt):} Licht, Ladung, Strom, Wärme
\item \textbf{Kl. 9 (30 USt):} Magnetismus, Induktion
\item \textbf{Kl. 10 (60 USt):} Beschleunigung, Wellen, Atom
\end{itemize}

\vspace{0.2cm}
\textbf{\color{physikblue}SEKUNDARSTUFE II}
\begin{itemize}[leftmargin=0.3cm, itemsep=1pt, topsep=2pt]
\item \textbf{Q1:} E-Feld, B-Feld, Induktion
\item \textbf{Q2:} Schwingungen, Wellen
\item \textbf{Q3:} Quantenphysik, Kernphysik
\item \textbf{Q4:} SRT (nur LK!)
\end{itemize}

\vspace{0.2cm}
\textbf{\color{physikblue}KOMPETENZEN}
\begin{itemize}[leftmargin=0.3cm, itemsep=1pt, topsep=2pt]
\item [S] Sachkompetenz
\item [E] Erkenntnisgewinnung
\item [K] Kommunikation
\item [B] Bewertung
\end{itemize}

\vspace{0.2cm}
\textbf{\color{physikblue}EXPERIMENTE}
\begin{itemize}[leftmargin=0.3cm, itemsep=1pt, topsep=2pt]
\item \SE{} = Schülerexperiment
\item \DE{} = Demonstrationsexperiment
\end{itemize}
\end{minipage}
\hfill
\begin{minipage}[t]{0.48\textwidth}
\textbf{\color{physikblue}AFB-VERTEILUNG ABITUR}
\begin{itemize}[leftmargin=0.3cm, itemsep=1pt, topsep=2pt]
\item AFB I: $\sim$30\% (nennen, beschreiben, berechnen)
\item AFB II: $\sim$40\% (erklären, auswerten, vergleichen)
\item AFB III: $\sim$30\% (begründen, beurteilen, entwickeln)
\end{itemize}

\vspace{0.2cm}
\textbf{\color{physikblue}WICHTIGE FORMELN}
\begin{align*}
F_G &= m \cdot g & (g = 10!) \\
E_{pot} &= m \cdot g \cdot h \\
E_{kin} &= \tfrac{1}{2} m v^2 \\
p &= m \cdot v \\
c &= \lambda \cdot f \\
E &= h \cdot f
\end{align*}

\vspace{0.2cm}
\textbf{\color{physikblue}KOPFRECHENBAR}
\begin{itemize}[leftmargin=0.3cm, itemsep=1pt, topsep=2pt]
\item $g = \SI{10}{\metre\per\second\squared}$ (nicht 9,81!)
\item 36 km/h = 10 m/s
\item 72 km/h = 20 m/s
\item $\sqrt{4}, \sqrt{9}, \sqrt{16}, \sqrt{25}...$
\end{itemize}

\vspace{0.2cm}
\textbf{\color{physikblue}KONTAKT}\\
IQ M-V, Fachbereich 4\\
\url{bildung-mv.de/rahmenplaene}
\end{minipage}

\vspace{0.5cm}
\begin{center}
\begin{tikzpicture}
\node[rectangle, rounded corners=3pt, fill=lightgray, draw=gray, minimum width=14cm, minimum height=0.8cm] {
  \small\textit{Erstellt auf Basis der Rahmenpläne Physik MV | Stand: Dezember 2025}
};
\end{tikzpicture}
\end{center}

\end{document}
